\begin{lemma}
\label{lemma-affineline-special}
Let $R$ be a ring. Let $f, g \in R[x]$ be polynomials.
Assume the leading coefficient of $g$ is a unit of $R$.
There exist elements $r_i\in R$, $i = 1\ldots, n$ such that
the image of $D(f) \cap V(g)$ in $\Spec(R)$ is
$\bigcup_{i = 1, \ldots, n} D(r_i)$.
\end{lemma}

\begin{proof}
Write $g = ux^d + a_{d-1}x^{d-1} + \ldots + a_0$, where
$d$ is the degree of $g$, and hence $u \in R^*$.
If $d=0$, then $g=u$ is a unit, so $V(g)$ and its image are empty;
the required union may be taken to have no terms. Suppose $d>0$ and
consider the original ring $A=R[x]/(g)$. It is finite free with basis
the images of $1,x,\ldots,x^{d-1}$. Indeed, if a polynomial has leading
term $bx^e$ with $e\geq d$, subtracting
$bu^{-1}x^{e-d}g$ cancels that term and lowers its degree. Iteration
gives a remainder of degree less than $d$, keeping $g$ and all its
coefficients unchanged. If a nonzero multiple $qg$ had degree less
than $d$, its leading coefficient would be the product of the nonzero
leading coefficient of $q$ with the unit $u$, hence nonzero, and its
degree would be $\deg(q)+d\geq d$, a contradiction. This proves both
spanning and independence. Consider multiplication by the image of
$f$ on $A$. This is an $R$-module map.
Hence we can let $P(T) \in R[T]$ be the characteristic polynomial
of this map. Write $P(T) = T^d + r_{d-1} T^{d-1} + \ldots + r_0$.
We claim that $r_0, \ldots, r_{d-1}$ have the desired property.
We will use below the property of characteristic polynomials
that
$$
\mathfrak p \in V(r_0, \ldots, r_{d-1})
\Leftrightarrow
\text{multiplication by }f\text{ is nilpotent on }
A \otimes_R \kappa(\mathfrak p).
$$
This was proved in Lemma \ref{lemma-characteristic-polynomial-prime}.

\medskip\noindent
Suppose $\mathfrak q\in D(f) \cap V(g)$, and let
$\mathfrak p = \mathfrak q \cap R$. Then there is a nonzero map
$A \otimes_R \kappa(\mathfrak p) \to \kappa(\mathfrak q)$ which
is compatible with multiplication by $f$.
And $f$ acts as a unit on $\kappa(\mathfrak q)$.
Thus we conclude $\mathfrak p \not \in  V(r_0, \ldots, r_{d-1})$.

\medskip\noindent
On the other hand, suppose that $r_i \not\in \mathfrak p$ for some
prime $\mathfrak p$ of $R$ and some $0 \leq i \leq d - 1$.
Then multiplication by $f$ is not nilpotent on the algebra
$A \otimes_R \kappa(\mathfrak p)$.
Hence there exists a prime ideal $\overline{\mathfrak q} \subset
A \otimes_R \kappa(\mathfrak p)$ not containing the image of $f$.
The inverse image of $\overline{\mathfrak q}$ in $R[x]$
is an element of $D(f) \cap V(g)$ mapping to $\mathfrak p$.
\end{proof}